\chapter{PET Scan}
\section*{What Is This Test?} Positron emission tomography (PET) uses a radioactive tracer to show tissue metabolism and function.\section*{Alternative Names} PET imaging; PET/CT; positron emission tomography.\section*{Principle} A tracer emits positrons detected by the scanner; CT or MRI may add anatomic detail.\section*{Why This Test Is Ordered} It helps stage or monitor cancer and evaluate selected heart, brain, and inflammatory disorders.\section*{Primary vs Secondary Test} PET is a secondary imaging test selected for a specific clinical question; it is not a general screening scan.\section*{How This Test Is Performed} Tracer is injected, the patient rests while it distributes, and images are acquired while lying still.\section*{What Preparation Is Needed} Fasting and glucose instructions are common. Report diabetes, pregnancy, breastfeeding, and medicines.\section*{Understanding Results} Uptake can reflect cancer, infection, inflammation, or normal physiology; findings require CT/MRI and clinical correlation.\section*{Follow-up Tests} Biopsy, diagnostic CT/MRI, laboratory tests, or oncology consultation may follow.\section*{References} \begin{enumerate}\item MedlinePlus. PET Scan.\item Society of Nuclear Medicine and Molecular Imaging. PET procedure standards.\end{enumerate}\clearpage\section*{Understanding Results and Follow-up}
The laboratory report should identify the specimen, method, units, reference interval or decision limit, and whether the result is positive, negative, abnormal, or inconclusive. Reference intervals vary with age, sex, pregnancy, specimen type, assay, and laboratory; a result outside a range is not automatically a diagnosis, and a result within a range does not always exclude disease. Discuss unexpected results with the ordering clinician, who can decide whether repeat or confirmatory testing, treatment, or referral is appropriate. Do not change medicines or delay urgent care based on an isolated result.
\section*{Regulatory Status and Authorization Date}This chapter describes a test category, not a specific commercial product. FDA, CDSCO, EU IVDR, and MHRA approval or registration dates therefore cannot be assigned without the assay manufacturer, product name, instrument, and intended use. For laboratory-developed tests, an individual agency approval date may not exist. See the Preface for the interpretation of regulatory status.
\clearpage
